% Audited insertion for AG.1, owned here for transfer to geometric_routes.
% No new claim number is needed.
\paragraph{Consequences of point separation and the finite spectral test.}
Once A1 has been proved, the unrestricted collision-dependent
implications in L5 hold as well: there is no pair of distinct
points satisfying their premise. In particular the unqualified
effective-divisor conclusion is no longer an open assertion.
The rationality qualifications above describe what the direct
correspondence argument proves without that theorem; they do
not weaken the now-established implication.

For clarity, the proposed spectral test has the following exact
meaning. Let $D/\mathbf F_q$ be smooth, proper and geometrically
connected of genus $h\ge1$, and write
\[
 P_D(T)=\det(1-TF\mid H^1_{\mathrm{et}}
       (D_{\overline{\mathbf F}_q},\mathbf Q_\nu)),
 \qquad
 H_D(X)=X^{2h}P_D(X^{-1}),\quad \nu\ne p.
\]
Then $J_D$ has a nonzero supersingular isogeny factor if and
only if one root $\alpha$ of $H_D$ has
$\alpha/\sqrt q$ a root of unity. Indeed, on a simple
isogeny factor this is exactly the supersingular
Newton-slope criterion; see
\cite[\S\S21.20--21.21(3)]{AC_Oort2008}. If one root has
this property, every conjugate in its simple factor does,
so that factor is supersingular. Conversely, a
supersingular factor supplies such roots.

This is a finite exact test: it is equivalent to
\[
 \deg\gcd\bigl(H_D(X),X^{2n}-q^n\bigr)>0
 \quad\hbox{for some }1\le n\le8h^2.
\]
To prove the bound, suppose $\beta=\alpha^2/q$ has exact
root-of-unity order $n$. Then
$\varphi(n)=[\mathbf Q(\beta):\mathbf Q]\le2h$.
For $n=\prod_\ell\ell^{a_\ell}$,
\[
 \frac{\varphi(n)^2}{n}
       =\prod_{\ell\mid n}\ell^{a_\ell-2}(\ell-1)^2
       \ge\frac12.
\]
Every odd-prime factor is at least one, and the factor
at two is at least $1/2$. Thus $n\le8h^2$.
Conversely, any nonconstant displayed greatest common
divisor gives $\alpha^{2n}=q^n$, which is precisely the
required root-of-unity condition.

For a constant abelian cover $D\to C$, the factorization
$P_D=\prod_\chi L(C,\chi,T)$ makes this the character
test in L5, with the trivial character included when
testing the whole Jacobian. Omitting it tests only the
new factor. A constant nonabelian cover requires all
irreducible Artin factors with their multiplicities;
a non-Galois cover requires its actual permutation
representation, or simply the polynomial $P_D$ itself.
The fixed-cover test supplies no existence theorem for
an etale cover having such a factor and no assertion
that the resulting maps jointly separate points.
Universal prime-to-$p$ supersingular correspondences
therefore remain the separate existence problem in A5.
